\documentclass[11pt]{article}
\usepackage[margin=1in]{geometry}
\usepackage{amsmath,amssymb,amsthm}
\usepackage[colorlinks=true,linkcolor=blue,urlcolor=blue,citecolor=blue]{hyperref}
\newtheorem{theorem}{Theorem}
\newtheorem{lemma}[theorem]{Lemma}
\newtheorem{corollary}[theorem]{Corollary}
\newtheorem{remark}[theorem]{Remark}
\newcommand{\tr}{\operatorname{tr}}
\newcommand{\rank}{\operatorname{rank}}
\newcommand{\sgn}{\operatorname{sgn}}
\newcommand{\op}{\mathrm{op}}
\newcommand{\R}{\mathbb R}
\newcommand{\C}{\mathbb C}
\newcommand{\HH}{\mathcal H}
\title{Small Schatten exponents for fractional matrix maps:\\matching bounds and a finite lower construction}
\author{Working research note}
\date{November 14, 2026}
\begin{document}
\maketitle
\begin{abstract}
Fix $0<\alpha<1$. For the fractional maps $f(x)=|x|^\alpha$ and
$f(x)=\sgn(x)|x|^\alpha$, we establish matching bounds of order $s^{-2}$
for the best constant in the power-$s$ Schatten H\"older inequality as
$s\downarrow0$. Equivalently, the usual H\"older constant has logarithm
$2s^{-1}\log(1/s)+O_\alpha(s^{-1})$. The lower bound already holds for
positive matrices differing by a rank-one positive matrix of norm one.
It follows from an explicit finite Toeplitz construction and an exact
Woodbury comparison with the Fr\'echet derivative. An elementary
imaginary-resolvent quadrature gives the upper bound on the whole real
line. The positive-matrix root-exponential upper bound has classical
antecedents in Birman--Koplienko--Solomyak (1975); it is not claimed as new.
Priority for the matching lower bound and the quantitative conclusion
remains under review.
\end{abstract}

\section{Statement and scope}
For a finite matrix $X$, write $\|X\|_s=(\tr|X|^s)^{1/s}$, including
the quasi-norm range $0<s<1$. Put
\[
 f_e(x)=|x|^\alpha,\qquad f_o(x)=\sgn(x)|x|^\alpha,
 \qquad f_o(0)=0.
\]
For $f=f_e,f_o$, define
\[
 \HH_{\alpha,s}(f)=\sup_{n\ge1}\sup_{A=A^*,\ B=B^*,\ A\ne B}
 \frac{\tr|f(A)-f(B)|^s}{\tr|A-B|^{\alpha s}}.       \tag{1}
\]
Let $\HH^+_{\alpha,s}$ denote the same supremum restricted to $A,B\succeq0$,
where both maps equal $x^\alpha$. The usual best norm constant is
$K_{\alpha,s}(f)=\HH_{\alpha,s}(f)^{1/s}$.

\begin{theorem}\label{thm:main}
For each fixed $0<\alpha<1$, there are positive constants $c_\alpha,C_\alpha$
such that, for every $0<s\le1$,
\[
 \frac{c_\alpha}{s^2}\le \HH^+_{\alpha,s}
 \le \HH_{\alpha,s}(f_e),\ \HH_{\alpha,s}(f_o)
 \le \frac{C_\alpha}{s^2}.                            \tag{2}
\]
In particular, for either map,
\[
 \log K_{\alpha,s}(f)
   =\frac{2}{s}\log\frac1s+O_\alpha\!\left(\frac1s\right).
                                                               \tag{3}
\]
The best power-$s$ constant for the rectangular matrix map
$X\mapsto X|X|^{\alpha-1}$ equals $\HH_{\alpha,s}(f_o)$ exactly.
At a zero singular value this map is defined to be zero.
\end{theorem}

The inequality with a finite unspecified constant is established work:
see Ricard~\cite{Ricard} and the broader functional-calculus results
of Huang--Sukochev--Zanin~\cite{HSZ} and McDonald--Sukochev~\cite{MS}.
For positive matrices, the root-exponential estimate behind the upper
bound follows from the truncated rank-one integral estimate in
\cite[Section~3, Lemma~3 and (3.10)]{BKS}. We give a self-contained proof to
make all parameter dependence visible. The finite lower construction
and its relation to older weighted Cauchy/Hankel asymptotics require
further priority assessment. This note asserts no improved matrix
multiplication exponent or solution of the general thin-product
relative low-rank approximation problem.

\section{Uniform compression of finite-rank differences}
All unlabelled constants in this section can depend on the fixed
$\alpha$ and on a fixed strip width, but not on the dimension or on
the spectra of the two matrices.

\begin{lemma}\label{lem:rank}
Let $A,B$ be finite Hermitian matrices, $E=A-B$, $a=\|E\|_\op>0$, and
$r=\rank(E)$. For either $f_e$ or $f_o$ and every integer $m\ge1$, there
is a matrix $C_m$ of rank at most $2r(2m+1)$ such that
\[
 \|f(A)-f(B)-C_m\|_\op
       \le C_\alpha e^{-c_\alpha\sqrt m}a^\alpha.    \tag{4}
\]
Consequently, with possibly different positive constants,
\[
 \sigma_j(f(A)-f(B))
       \le C_\alpha e^{-c_\alpha\sqrt{j/r}}a^\alpha,
       \qquad j\ge1.                                \tag{5}
\]
\end{lemma}
\begin{proof}
For real $x$, set
\[
 F_t(x)=\frac{x^2}{x^2+t^2},\qquad
 G_t(x)=\frac{tx}{x^2+t^2}.
\]
The scalar integral identities are
\begin{align}
 |x|^\alpha
    &=\frac{2\sin(\pi\alpha/2)}{\pi}
          \int_0^\infty t^{\alpha-1}F_t(x)\,dt,\\
 \sgn(x)|x|^\alpha
    &=\frac{2\cos(\pi\alpha/2)}{\pi}
          \int_0^\infty t^{\alpha-1}G_t(x)\,dt.       \tag{6}
\end{align}
They follow by substituting $t=|x|u$ and the elementary beta integral;
the spectral theorem gives the corresponding matrix identities.

Write $R_A(z)=(A-zI)^{-1}$. The resolvent identity gives
\[
 R_A(z)-R_B(z)=-R_A(z)E R_B(z),                       \tag{7}
\]
whose rank is at most $r$. Moreover,
\begin{align}
 F_t(A)-F_t(B)
    &=-\frac{t}{2i}\{[R_A(it)-R_B(it)]-[R_A(-it)-R_B(-it)]\},\\
 G_t(A)-G_t(B)
    &=\frac{t}{2}\{[R_A(it)-R_B(it)]+[R_A(-it)-R_B(-it)]\}.
                                                               \tag{8}
\end{align}
Each difference therefore has rank at most $2r$.

Fix $0<d<\pi/2$ and allow $t=\rho e^{i\theta}$ with $|\theta|\le d$.
For real $x$,
\[
 |x^2+t^2|\ge \cos(d)(x^2+\rho^2).                 \tag{9}
\]
Indeed the square of the left side equals
$(x^2+\rho^2)^2-4x^2\rho^2\sin^2\theta$.
Thus $\|F_t(A)\|\le1/\cos d$ and
$\|G_t(A)\|\le1/(2\cos d)$. On the other hand,
$\|R_A(\pm it)\|\le1/(\rho\cos d)$, so (7)--(8) give
\[
 \|F_t(A)-F_t(B)\|,\ \|G_t(A)-G_t(B)\|
       \le C_d\min(1,a/\rho).                       \tag{10}
\]
No norm of $A$ or $B$ occurs.

Substitute $t=a e^y$ in the difference of either identity (6). Its
integrand $\mathcal F(y)$ is analytic on $|\Im y|<\pi/2$, and (10) yields
\[
 \|\mathcal F(x+i\theta)\|
   \le C_{\alpha,d}a^\alpha
       \min(e^{\alpha x},e^{-(1-\alpha)x}),
       \qquad |\theta|\le d.                        \tag{11}
\]
In particular its boundary integrals are bounded by
$C_{\alpha,d}a^\alpha$.

For completeness, the strip trapezoidal estimate used here is
\[
 \left\|\int_\R\mathcal F(x)\,dx
               -h\sum_{j\in\mathbb Z}\mathcal F(jh)\right\|
   \le \frac{2M}{e^{2\pi d/h}-1},                   \tag{12}
\]
when both boundary integrals are at most $M$ and the integrand has
integrable exponential tails. Shift the Fourier transform contour to
$\Im y=\pm d$ to bound its transform by $M e^{-d|\xi|}$, and apply
Poisson summation. This proves (12), also for matrices: the same
dimension-independent bounds apply after pairing with any two unit vectors.

Take $h=m^{-1/2}$ and truncate the sum to $-m\le j\le m$.
Equation (11) bounds the omitted tails by
$C_\alpha a^\alpha e^{-c_\alpha\sqrt m}$, while (12) gives the same
form for the infinite-sum error. Each retained summand has rank at
most $2r$, proving (4). The best-rank spectral approximation theorem
then gives (5) after adjusting constants. For the finitely many initial
indices, integrate (11) directly to obtain
$\|f(A)-f(B)\|\le C_\alpha a^\alpha$.
\end{proof}

\begin{remark}
For $A,B\succeq0$ and $f(x)=x^\alpha$, the same conclusion is obtained
with negative real poles and rank-$r$ summands from
\[
 B^\alpha-A^\alpha
   =\frac{\sin(\pi\alpha)}{\pi}\int_0^\infty
      t^\alpha[(A+tI)^{-1}-(B+tI)^{-1}]\,dt.
                                                               \tag{13}
\]
This also makes the classical connection to \cite{BKS} explicit.
\end{remark}

\section{From rank one to the upper H\"older constant}
For a rank-one Hermitian perturbation $E$ of norm $a$, (5) and monotonicity
of $x\mapsto e^{-c s\sqrt x}$ imply, for $0<s\le1$,
\begin{align}
 \tr|f(A)-f(B)|^s
 &\le C_\alpha^s a^{\alpha s}\sum_{j\ge1}e^{-c_\alpha s\sqrt j}\\
 &\le C_\alpha^s a^{\alpha s}
                  \int_0^\infty e^{-c_\alpha s\sqrt x}\,dx
 \le \frac{U_\alpha}{s^2}a^{\alpha s}.              \tag{14}
\end{align}
The last integral is $2/(c_\alpha^2s^2)$; replacing $C_\alpha^s$ by
$\max(1,C_\alpha)$ leaves $U_\alpha$ independent of $s$.

For general Hermitian $E=A-B$, take its spectral decomposition
$E=\sum_{j=1}^r t_j u_j u_j^*$ and let
$A_j=B+\sum_{\ell\le j}t_\ell u_\ell u_\ell^*$.
The Schatten power triangle inequality (Rotfel'd's inequality) is
\[
 \tr|X+Y|^s\le\tr|X|^s+\tr|Y|^s,\qquad0<s\le1.
\]
Apply (14) to each difference $f(A_j)-f(A_{j-1})$ and telescope:
\[
 \tr|f(A)-f(B)|^s
 \le\frac{U_\alpha}{s^2}\sum_j|t_j|^{\alpha s}
 =\frac{U_\alpha}{s^2}\tr|A-B|^{\alpha s}.          \tag{15}
\]
No factor depending on the number of summands is introduced.
This proves the upper bound in Theorem~\ref{thm:main}.
For a proof confined to positive matrices, the positive and negative
parts of $E$ can instead be telescoped along the two positive paths
through the common upper endpoint $B+E_+=A+E_-$.

\section{An exact finite lower construction}
The next result is sufficient for every lower bound in the main theorem.
\begin{theorem}\label{thm:finite-lower}
Fix $0<\alpha<1$, write $\beta=1-\alpha$, and let $k\ge1$ be an integer.
There are a positive diagonal $k$-by-$k$ matrix $H\succeq I$ and a
rank-one positive matrix $E$ with $\|E\|_\op=1$ such that
\[
 \lambda_{\min}((H+E)^\alpha-H^\alpha)
 \ge \frac{\pi\beta\sin(\pi\alpha)}{2}
          \exp\{-2\pi\sqrt{2\beta k}\}.             \tag{16}
\]
\end{theorem}
\begin{proof}
We first give an exact comparison that avoids a perturbative remainder.
If $E=vv^*$, $\|v\|=1$, the Woodbury identity in (13) gives
\[
 (H+E)^\alpha-H^\alpha
 =\frac{\sin(\pi\alpha)}{\pi}\int_0^\infty
 \frac{t^\alpha(H+tI)^{-1}vv^*(H+tI)^{-1}}
      {1+v^*(H+tI)^{-1}v}\,dt.
\]
Because $H\succeq I$, the denominator lies in $[1,2]$. The same
integral without its denominator is the Fr\'echet derivative
$D(H^\alpha)[E]$. Consequently
\[
 (H+E)^\alpha-H^\alpha\succeq \tfrac12 D(H^\alpha)[E]. \tag{17}
\]

Temporarily let $h>0$ be arbitrary. For $0\le i<k$, set
\[
 \lambda_i=e^{hi},\qquad H=\operatorname{diag}(\lambda_i),
 \qquad w_i=\alpha^{-1/2}\lambda_i^{\beta/2},
 \qquad N=\|w\|^2,\qquad v=w/\sqrt N.
\]
The divided-difference formula for the derivative yields
\[
 D(H^\alpha)[E]=N^{-1}T,\qquad
 T_{ij}=\kappa_\alpha(h(i-j)),\qquad
 \kappa_\alpha(x)=\frac{\sinh(\alpha x/2)}
                       {\alpha\sinh(x/2)},\quad\kappa_\alpha(0)=1.
                                                               \tag{18}
\]
For example, on the diagonal this is
$\alpha\lambda_i^{\alpha-1}|v_i|^2=1/N$.

With the Fourier convention
$\widehat\kappa(\xi)=\int_\R e^{-i\xi x}\kappa(x)\,dx$, one has
\[
 \widehat\kappa_\alpha(\xi)
 =\frac{2\pi\sin(\pi\alpha)}
        {\alpha[\cosh(2\pi\xi)+\cos(\pi\alpha)]}
 \ge\frac{\pi\sin(\pi\alpha)}{\alpha}e^{-2\pi|\xi|}.
                                                               \tag{19}
\]
The identity follows by expanding, for $x>0$,
\[
 \kappa_\alpha(x)=\alpha^{-1}\sum_{n\ge0}
 \left(e^{-(n+\beta/2)x}-e^{-(n+(1+\alpha)/2)x}\right)
\]
and evaluating the resulting partial fractions with the cotangent
identity; it is also the Fourier formula for the strip Poisson kernel.
The formula gives the check
$\kappa_{1/2}(x)=\operatorname{sech}(x/4)$ and
$\widehat\kappa_{1/2}(\xi)=4\pi\operatorname{sech}(2\pi\xi)$.

Poisson summation gives the positive Toeplitz symbol
\[
 p_h(\theta)=\sum_{j\in\mathbb Z}\kappa_\alpha(hj)e^{ij\theta}
   =h^{-1}\sum_{\ell\in\mathbb Z}
        \widehat\kappa_\alpha((\theta+2\pi\ell)/h).
\]
For $|\theta|\le\pi$, its $\ell=0$ term and (19) imply
\[
 p_h(\theta)\ge\frac{\pi\sin(\pi\alpha)}{\alpha h}
                    e^{-2\pi^2/h}.
\]
The quadratic form of the \emph{finite} $k$-by-$k$ Toeplitz matrix
$T$ is the integral of $p_h$ times a trigonometric polynomial's squared
modulus. Parseval therefore proves
\[
 \lambda_{\min}(T)\ge\frac{\pi\sin(\pi\alpha)}{\alpha h}
                    e^{-2\pi^2/h}.                 \tag{20}
\]
No infinite-section eigenvalue asymptotics or boundary correction is used.

Finally,
\[
 N=\alpha^{-1}\sum_{i=0}^{k-1}e^{\beta hi}
   \le\frac{e^{\beta hk}}{\alpha\beta h}.
\]
Combining this with (17)--(20) gives
\[
 \lambda_{\min}((H+E)^\alpha-H^\alpha)
 \ge\frac{\pi\beta\sin(\pi\alpha)}2
           e^{-2\pi^2/h-\beta hk}.
\]
Choose $h=\pi\sqrt{2/(\beta k)}$ to obtain (16).
\end{proof}

Write $b_\alpha=\pi(1-\alpha)\sin(\pi\alpha)/2$ and
$d_\alpha=2\pi\sqrt{2(1-\alpha)}$. Since $E$ has norm and rank one,
Theorem~\ref{thm:finite-lower} implies
\[
 \HH^+_{\alpha,s}\ge k b_\alpha^s e^{-d_\alpha s\sqrt k}. \tag{21}
\]
Taking $k=\lceil s^{-2}\rceil$ proves the lower bound in (2).
More precisely, choosing integers $k\sim4/(d_\alpha^2s^2)$ gives
\[
 \liminf_{s\downarrow0}s^2\HH^+_{\alpha,s}
       \ge\frac{1}{2\pi^2(1-\alpha)e^2}.            \tag{22}
\]
The exact leading constant is not established here.

\begin{corollary}
For every fixed $0<\alpha<1$, the worst-case approximation rank for
$B^\alpha-A^\alpha$ at operator error
$\varepsilon\|B-A\|_\op^\alpha$, when $A,B\succeq0$ and
$\rank(B-A)\le r$, is of order
$r\log^2(1/\varepsilon)$ as $\varepsilon\downarrow0$, until limited
by the ambient dimension. The lower bound already holds for positive updates.
\end{corollary}
\begin{proof}
The positive version of Lemma~\ref{lem:rank} gives the upper bound.
For $r=1$, choose $k$ a sufficiently small positive multiple of
$\log^2(1/\varepsilon)$ in (16). All $k$ singular values then exceed
$\varepsilon$, while $\|E\|=1$. The spectral best-rank theorem forces
rank at least $k$. Direct sums of $r$ identical copies preserve the
operator norm of the update and give the additional factor $r$.
\end{proof}

\section{Rectangular matrices and limitations}
For rectangular $X$, use the Hermitian dilation
\[
 \mathcal D(X)=\begin{pmatrix}0&X\\X^*&0\end{pmatrix}.
\]
Then $f_o(\mathcal D(X))=\mathcal D(X|X|^{\alpha-1})$.
Every nonzero singular value is repeated twice by dilation, so
\[
 \tr|\mathcal D(X)-\mathcal D(Y)|^{\alpha s}
       =2\tr|X-Y|^{\alpha s},
\]
and the identical factor $2$ occurs in the numerator. This proves that
the rectangular constant is at most the Hermitian odd-map constant.
The reverse inequality follows by taking square Hermitian inputs.

The estimates concern the best constant as $s\downarrow0$ with
$\alpha$ fixed. They give its sharp power order and the leading
coefficient $2$ in (3), not the exact multiplicative constant in
$s^2\HH_{\alpha,s}$. Finite Schatten matrices suffice for the lower
bound. A semifinite-algebra extension can be formulated using trace
of range projections in place of rank; type III algebras need additional
arguments and are not included in the theorem as stated.

The compression norm is
$\varepsilon\|B-A\|^\alpha$, rather than
$\varepsilon\|B^\alpha-A^\alpha\|$. A rank-one update can have a
full relative-rank correction. Nor does a quadrature rank bound alone
pay for the shifted linear solves required to construct its factors.
The separate raw-covariance algorithmic notes give conventional-arithmetic
implementations for a positive square-root update, with all such solves
charged. Neither statement solves the nonzero real-threshold problem
for arbitrary matrix-product low-rank approximation.

\begin{thebibliography}{9}
\bibitem{BKS}
M. Sh. Birman, L. S. Koplienko, M. Z. Solomyak.
Estimates of the spectrum of a difference of fractional powers of
selfadjoint operators. Izv. Vyssh. Uchebn. Zaved. Mat., no.~3, 3--10, 1975.
\url{https://www.mathnet.ru/eng/ivm6330}.
\bibitem{Ricard}
\'E. Ricard. Fractional powers on noncommutative $L_p$ for $p<1$.
Adv. Math. 333, 194--211, 2018.
\url{https://arxiv.org/abs/1805.05677}.
\bibitem{HSZ}
J. Huang, F. Sukochev, D. Zanin.
Operator $\theta$-H\"older functions with respect to $\|\cdot\|_p$,
$0<p\le\infty$. J. London Math. Soc., 2022.
\url{https://arxiv.org/abs/2110.09708}.
\bibitem{MS}
E. McDonald, F. Sukochev. Lipschitz estimates in quasi-Banach
Schatten ideals. Math. Ann., 2022.
\url{https://arxiv.org/abs/2009.08069}.
\end{thebibliography}
\end{document}
